\documentclass[11pt]{article}

\usepackage[utf8]{inputenc}
\usepackage[T1]{fontenc}
\usepackage{lmodern}
\usepackage[margin=1in]{geometry}
\usepackage{setspace}
\usepackage{graphicx}
\usepackage{booktabs}
\usepackage{array}
\usepackage{float}
\usepackage{hyperref}
\usepackage{xcolor}
\usepackage{enumitem}
\usepackage{amsmath}
\usepackage{amssymb}
\usepackage{amsthm}

\hypersetup{
  colorlinks=true,
  linkcolor=blue!50!black,
  citecolor=blue!50!black,
  urlcolor=blue!50!black,
  pdftitle={Recurrence Diagnostics for Econometric Stress Episodes}
}

\setstretch{1.08}
\setlength{\parskip}{0.55em}
\setlength{\parindent}{0pt}
\emergencystretch=2em

\newtheorem{proposition}{Proposition}
\newtheorem{remark}{Remark}

\newcommand{\keywords}[1]{\noindent\textbf{Keywords:} #1}
\newcommand{\figasset}[2]{%
  \begin{figure}[H]
    \centering
    \IfFileExists{#1}{%
      \includegraphics[width=0.92\textwidth]{#1}%
    }{%
      \fbox{%
        \parbox[c][0.24\textheight][c]{0.86\textwidth}{%
          \centering
          Missing figure asset.\\
          Expected file: \texttt{\detokenize{#1}}%
        }%
      }%
    }
    \caption{#2}
  \end{figure}
}

\title{Recurrence Diagnostics for Econometric Stress Episodes:\\
A Methods Paper with a Cluster-Adjusted Stress Functional}
\author{Author Name\\
Institution\\
\texttt{author@example.com}}
\date{July 2026}

\begin{document}

\maketitle

\begin{abstract}
Standard stress diagnostics in applied econometrics typically measure the magnitude of variation, tail losses, or event frequency, but not the temporal organization of extreme episodes. That omission matters whenever policy or risk interpretation depends on whether extremes arrive as isolated shocks or as clustered stress. This paper develops a reproducible recurrence-oriented framework for stress diagnostics and introduces a cluster-adjusted stress functional that combines exceedance probability with tail clustering. The paper has two goals. First, it provides a serious methods contribution: a formal probabilistic setup, explicit estimators for exceedance incidence, return-time behavior, runs-based clustering, and multivariate stress activation, and a reproducible implementation in the accompanying repository. Second, it adds a narrow theoretical contribution: basic ordering and plug-in consistency results for the proposed cluster-adjusted functional under standard high-threshold regularity conditions. Synthetic evidence generated by the repository shows that recurrence diagnostics separate benchmark processes that would be inadequately summarized by volatility alone. In the current artifact set, the estimated extremal indices for two stress benchmarks are $0.673$ and $0.780$, the return-time distribution has median $14$ and 90th percentile $40.2$, and rolling volatility is elevated during exceedance periods without exhausting the clustering information. The paper makes modest claims: it does not offer new asymptotic theory for extremal-index estimation itself, nor does it claim structural macroeconomic identification. Its contribution is a mathematically explicit and computationally reproducible bridge between extreme-value diagnostics and econometric stress analysis.
\end{abstract}

\keywords{extreme values, extremal index, recurrence, return times, stress testing, time series econometrics}

\section{Introduction}

The motivating problem is straightforward. In many macro-financial and econometric applications, a ``stress episode'' is treated as a period in which a variable becomes large in magnitude, volatile, or both. But a stress episode is not defined only by how large observations become. It is also defined by how they recur. A series that produces isolated tail observations conveys a different kind of fragility from a series that produces clustered extremes, even if the two series have comparable event counts or similar rolling volatility.

That distinction is conceptually obvious and methodologically underused. Standard variance-based diagnostics, from ARCH and GARCH models onward, are designed to summarize time-varying second moments \cite{engle1982,bollerslev1986}. Tail-risk and systemic-risk measures such as VaR, CoVaR, expected shortfall, and capital-shortfall metrics focus on the severity of bad states and their cross-sectional comovement \cite{adrian2016,acharya2017,brownlees2017,bisias2012}.

Yet the temporal spacing of threshold exceedances is a different object. Extreme-value theory for dependent sequences treats exceedance clustering and return times as first-order features of the data-generating process rather than as narrative afterthoughts \cite{leadbetter1983,embrechts1997,coles2001}. The extremal index provides a canonical summary of tail clustering \cite{obrien1987,ferro2003,suveges2007}. Related recurrence ideas also appear in the study of hitting times and rare events for stochastic and dynamical systems \cite{hirata1999,freitas2010,lucarini2016}. What is comparatively rare in applied econometrics is a workflow that makes these objects operational in a transparent, reproducible way and places them next to familiar volatility baselines rather than outside the analysis.

This paper addresses that gap. Its first contribution is methodological. I formalize a recurrence-oriented stress framework built around five objects: threshold exceedance incidence, return-time distributions, runs-based tail clustering, multivariate stress activation, and baseline volatility diagnostics. The framework is implemented in the accompanying repository with deterministic figure and CSV generation. Its second contribution is theoretical but deliberately narrow. I introduce a cluster-adjusted stress functional that scales exceedance incidence by tail clustering, prove a basic monotonicity result, and establish plug-in consistency under standard regularity conditions. These results are not a replacement for the deep theory of dependent extremes; they are a bridge from that theory to a usable stress-analysis summary.

The paper's empirical claims remain modest. The current repository supports synthetic evidence and a small multivariate illustration, not a structural macroeconomic application. The scientific claim is therefore limited to the one the evidence can bear: recurrence diagnostics recover dimensions of stress persistence and organization that are not exhausted by volatility summaries alone.

The remainder of the paper is organized as follows. Section 2 reviews the relevant literature and identifies the specific gap addressed here. Section 3 introduces the probabilistic framework and the estimands. Section 4 describes the estimators and the implementation logic. Section 5 develops the cluster-adjusted stress functional and states basic theoretical results. Section 6 presents synthetic evidence from the generated artifacts. Section 7 discusses limitations and scope.

\section{Literature Review}

\subsection{Extreme values under dependence}

The first literature strand is extreme-value theory for dependent sequences. Classical EVT is often introduced through independent observations, but economically meaningful time series are rarely independent. Leadbetter, Lindgren, and Rootz\'en \cite{leadbetter1983} established the modern framework for extremes of dependent sequences, including the role of mixing conditions and clustering. Embrechts, Kl\"uppelberg, and Mikosch \cite{embrechts1997} and Coles \cite{coles2001} made these ideas broadly accessible for finance, insurance, and applied risk analysis. The core lesson for the present paper is that dependence changes the interpretation of tail events: once exceedances cluster, the effective informational content of an event count changes as well.

\subsection{Extremal index and clustering diagnostics}

The second strand concerns the extremal index and practical cluster diagnostics. O'Brien \cite{obrien1987} formalized the extremal index as a scalar descriptor of clustering intensity for stationary sequences. Ferro and Segers \cite{ferro2003} and S\"uveges \cite{suveges2007} developed influential estimation strategies for dependent exceedances. The theoretical object is now standard, but its applied use in econometric stress workflows remains uneven. Many empirical analyses still use threshold exceedance counts, peak-over-threshold magnitudes, or volatility episodes without explicitly quantifying clustering. The present paper borrows from this literature but does not attempt to improve extremal-index estimation itself. Instead, it uses runs-based clustering summaries as ingredients in a stress diagnostic designed for applied interpretability.

\subsection{Recurrence and return-time analysis}

The third strand studies recurrence directly. In stochastic-process and dynamical-systems settings, return times, hitting times, and rare-event recurrence have long been treated as structurally meaningful objects. Hirata, Saussol, and Vaienti \cite{hirata1999} linked return-time statistics to recurrence properties in dynamical systems. Freitas, Freitas, and Todd \cite{freitas2010} and Lucarini et al. \cite{lucarini2016} deepened the connection between recurrence, extremes, and temporal dependence. This literature is mathematically rich, but much of it remains separated from mainstream econometric stress analysis. The separation is partly cultural and partly computational: applied pipelines rarely store or visualize recurrence summaries as first-class outputs.

\subsection{Volatility econometrics and tail-risk baselines}

The fourth strand is the dominant econometric baseline. ARCH and GARCH models \cite{engle1982,bollerslev1986} made time-varying volatility central to empirical finance and macroeconometrics. Later risk measures, including CoVaR \cite{adrian2016}, translate tail severity and conditional downside dependence into policy-relevant objects. These methods are indispensable, and the current paper does not argue against them. The issue is narrower: scale and conditional tail-loss measures are not designed to encode the spacing and clustering of repeated threshold exceedances. A period of elevated volatility can contain either isolated shocks or persistent recurrence. If stress persistence matters, recurrence should be measured rather than verbally inferred.

\subsection{Macro-financial stress monitoring and systemic-risk measurement}

The fifth strand concerns the broader literature on stress monitoring and systemic risk. Surveys such as Bisias et al. \cite{bisias2012} show that systemic-risk measurement already employs a wide menu of market-based, balance-sheet-based, and network-based summaries. Measures such as CoVaR \cite{adrian2016}, SRISK \cite{brownlees2017}, and systemic expected shortfall \cite{acharya2017} are designed to quantify the severity of financial distress and its spillovers across institutions. These are important objects, but they mainly answer questions about exposure, capital shortfall, and cross-sectional dependence. They do not directly summarize how often stress reappears once it has begun or how tightly grouped those recurrences are through time. The present paper should therefore be read as complementary to systemic-risk monitoring, not as a substitute for it.

\subsection{Gap and contribution}

The gap is therefore not an absence of theory about extremes, nor an absence of econometric tools for volatility and tail risk. The gap is the lack of a disciplined, reproducible bridge between those domains for the specific task of stress-episode characterization. This paper contributes to that bridge in three ways:
\begin{enumerate}[leftmargin=1.4em]
  \item it defines a recurrence-oriented stress framework in explicitly probabilistic terms;
  \item it implements that framework in repository-native code that generates article-facing evidence;
  \item it introduces a cluster-adjusted stress functional that converts clustering information into a directly interpretable stress measure.
\end{enumerate}

\section{Probabilistic Framework}

Let $\{X_t\}_{t \in \mathbb{Z}}$ be a strictly stationary real-valued process on a probability space $(\Omega,\mathcal{F},\mathbb{P})$ with marginal distribution function $F$. The object $X_t$ may represent a transformed return, a stress indicator, or another scalar measure whose upper tail corresponds to more severe conditions. The multivariate extension is introduced below.

For a threshold $u$, define the exceedance indicator
\[
I_t(u)=\mathbf{1}\{X_t>u\}.
\]
The exceedance probability is
\[
p(u)=\mathbb{P}(X_0>u).
\]
In finite samples of length $n$, high-threshold analysis typically works with a threshold sequence $u_n$ satisfying
\[
n\,\overline{F}(u_n)\to \tau \in (0,\infty),
\]
where $\overline{F}(u)=1-F(u)$. This keeps exceedances rare but observable.

If $t_1<\cdots<t_{N_n}$ denote the exceedance times among $\{1,\dots,n\}$, the return times are
\[
\tau_{j,n}=t_{j+1}-t_j,\qquad j=1,\dots,N_n-1.
\]
These random intervals summarize recurrence to the stress region. Their empirical distribution captures spacing, dispersion, and asymmetry in a way that event counts cannot.

To describe clustering, fix a run parameter $r\geq 1$. A new cluster begins at time $t$ if $X_t>u_n$ and the previous $r$ observations do not exceed $u_n$. The corresponding finite-level cluster-start indicator is
\[
C_{t,n}^{(r)}=I_t(u_n)\prod_{j=1}^{r}\left(1-I_{t-j}(u_n)\right).
\]
The associated finite-level clustering coefficient is
\[
\theta_n^{(r)}=\frac{\mathbb{E}[C_{0,n}^{(r)}]}{\mathbb{E}[I_0(u_n)]}
=\mathbb{P}\!\left(X_{-1}\le u_n,\dots,X_{-r}\le u_n \,\middle|\, X_0>u_n\right).
\]
When the standard runs approximation is appropriate, $\theta_n^{(r)}$ is a finite-level analogue of the extremal index: lower values indicate stronger local clustering of exceedances.

For a $d$-variate process $\{X_{t,j}\}$ with component-specific thresholds $u_{j,n}$, define
\[
I_{t,j}(u_{j,n})=\mathbf{1}\{X_{t,j}>u_{j,n}\},
\qquad
S_t=\frac{1}{d}\sum_{j=1}^{d} I_{t,j}(u_{j,n}).
\]
The scalar $S_t \in [0,1]$ is a multivariate stress-state index: it records the share of active stress coordinates at time $t$. It is descriptive, not causal.

\section{Estimators and Implementation}

\subsection{Sample estimators}

Given observations $X_1,\dots,X_n$, the empirical exceedance rate at threshold $u_n$ is
\[
\widehat{p}_n(u_n)=\frac{1}{n}\sum_{t=1}^{n}I_t(u_n).
\]
The runs-based estimator of the finite-level clustering coefficient is
\[
\widehat{\theta}_n^{(r)}=
\frac{\sum_{t=r+1}^{n} C_{t,n}^{(r)}}{\sum_{t=1}^{n} I_t(u_n)},
\]
whenever the denominator is positive. This estimator is simple, transparent, and computationally appropriate for a reproducible pipeline. It is not always the most statistically efficient choice, but it is easy to audit and aligns with the article's emphasis on inspectable artifacts.

The empirical return-time distribution is
\[
\widehat{G}_n(x)=\frac{1}{N_n-1}\sum_{j=1}^{N_n-1}\mathbf{1}\{\tau_{j,n}\le x\},
\]
provided $N_n\ge 2$. Summary functionals such as the mean, median, upper quantiles, and maximum of $\{\tau_{j,n}\}$ are used throughout the evidence section.

For the multivariate setting, the empirical stress-state index is
\[
\widehat{S}_t=\frac{1}{d}\sum_{j=1}^{d}I_{t,j}(u_{j,n}),
\]
and joint stress is defined by the event $\{\widehat{S}_t \geq 2/d\}$ in the repository's current illustration. This threshold can be varied; the key point is to make the rule explicit.

\subsection{Threshold and run-length choices}

Threshold selection is a substantive modeling decision, not a clerical preprocessing step. A threshold set too low destroys the rare-event interpretation; a threshold set too high produces severe finite-sample instability. The repository therefore treats threshold sensitivity as part of the diagnostic output rather than as a hidden implementation detail. In applied work, one should report how clustering summaries move across a threshold grid and how the chosen run length $r$ affects cluster identification.

The run parameter $r$ has an equally direct interpretation. Larger values merge nearby exceedances into common clusters and therefore tend to increase measured persistence. This does not make large $r$ ``better.'' It means that $r$ encodes what the analyst wishes to count as persistence. A serious stress paper should disclose that choice and test its sensitivity.

\subsection{Repository implementation}

The repository generates figure and CSV artifacts through a deterministic script that proceeds from raw or simulated series to threshold definition, recurrence summaries, baseline econometric diagnostics, and article-facing outputs. The conceptual pipeline contains eight recorded stages: raw series, transformations, stress-region definition, exceedance events, clustering diagnostics, return-time diagnostics, econometric baselines, and article evidence. This separation matters because recurrence statistics are sensitive to thresholding and clustering rules; saved intermediate tables are therefore part of the method.

\figasset{figures/01_conceptual_pipeline.png}{Conceptual pipeline for the repository-native workflow. The methodological point is that recurrence objects are treated as explicit data products rather than informal visual impressions.}

\section{A Cluster-Adjusted Stress Functional}

\subsection{Definition}

The paper's narrow original contribution is the following finite-level stress functional:
\[
\Lambda(u_n,r)=\frac{p(u_n)}{\theta_n^{(r)}}.
\]
Its sample analogue is
\[
\widehat{\Lambda}_n(u_n,r)=\frac{\widehat{p}_n(u_n)}{\widehat{\theta}_n^{(r)}},
\]
whenever $\widehat{\theta}_n^{(r)}>0$.

The interpretation is direct. The exceedance rate $\widehat{p}_n$ measures how often the process enters the stress region. Dividing by $\widehat{\theta}_n^{(r)}$ inflates that rate when exceedances are more clustered. Since $\theta_n^{(r)} \le 1$, the functional satisfies $\Lambda(u_n,r)\ge p(u_n)$. It may be read as a cluster-adjusted stress load: the same incidence rate is treated as more severe when it is delivered in more persistent bursts.

This definition is intentionally simple. It is not meant to replace richer tail-process theory. It is meant to convert a dependence object into an interpretable scalar that can sit next to frequency and volatility summaries.

\subsection{Basic properties}

\begin{proposition}[Monotone ordering by clustering]
Fix a threshold $u$ and run length $r$. Consider two stationary processes $A$ and $B$ such that
\[
p_A(u)=p_B(u)=p(u)>0,
\qquad
\theta_A^{(r)}(u)<\theta_B^{(r)}(u)\le 1.
\]
Then
\[
\Lambda_A(u,r)>\Lambda_B(u,r).
\]
\end{proposition}

\textit{Proof sketch.}
The conclusion follows immediately from the definition $\Lambda(u,r)=p(u)/\theta^{(r)}(u)$ and positivity of $p(u)$. For equal exceedance probability, the process with smaller clustering coefficient has the larger cluster-adjusted stress load. The content of the proposition is interpretive rather than algebraically deep: equal event frequency should not imply equal stress severity when one process exhibits stronger local bunching of extremes.

\begin{proposition}[Independent benchmark]
If exceedances are temporally independent at threshold $u$, then for every $r\ge 1$,
\[
\theta^{(r)}(u)=1
\qquad\text{and}\qquad
\Lambda(u,r)=p(u).
\]
\end{proposition}

\textit{Proof sketch.}
Under independence,
\[
\mathbb{P}\!\left(X_{-1}\le u,\dots,X_{-r}\le u \mid X_0>u\right)=1,
\]
so the finite-level clustering coefficient equals one. The cluster adjustment then disappears. This proposition supplies the correct benchmark: the functional adds information only when temporal dependence in extremes is present.

\begin{proposition}[Plug-in consistency]
Assume:
\begin{enumerate}[label=(A\arabic*),leftmargin=1.6em]
  \item $\{X_t\}$ is strictly stationary and strongly mixing with coefficients sufficient for a law of large numbers for the exceedance and cluster-start indicators;
  \item $u_n$ is a high-threshold sequence such that $p(u_n)>0$ for all $n$ and $p(u_n)\to 0$;
  \item for a fixed run length $r$, $\theta_n^{(r)}\to \theta^{(r)} \in (0,1]$;
  \item $\widehat{p}_n(u_n)\xrightarrow{p} p(u_n)$ and $\widehat{\theta}_n^{(r)}\xrightarrow{p}\theta^{(r)}$.
\end{enumerate}
Then
\[
\widehat{\Lambda}_n(u_n,r)-\frac{p(u_n)}{\theta^{(r)}}\xrightarrow{p}0.
\]
If, in addition, $p(u_n)\to p^\star$ and $\theta^{(r)}\to \theta^\star>0$, then
\[
\widehat{\Lambda}_n(u_n,r)\xrightarrow{p}\frac{p^\star}{\theta^\star}.
\]
\end{proposition}

\textit{Proof sketch.}
Under the stated assumptions, the numerator and denominator converge in probability to positive limits. The continuous mapping theorem applied to $(x,y)\mapsto x/y$ on $\mathbb{R}\times(0,\infty)$ gives the result. The proposition is intentionally conservative: it does not claim new asymptotic theory for $\widehat{\theta}_n^{(r)}$, only that once consistent ingredients are available, the plug-in stress functional is itself consistent.

\begin{proposition}[Run-length monotonicity]
Fix a threshold $u_n$ and two run lengths $1 \le r_1 < r_2$. Then
\[
\theta_n^{(r_2)} \le \theta_n^{(r_1)}.
\]
Consequently, whenever both quantities are positive,
\[
\Lambda(u_n,r_2) \ge \Lambda(u_n,r_1).
\]
\end{proposition}

\textit{Proof sketch.}
By definition,
\[
\theta_n^{(r)}=
\mathbb{P}\!\left(X_{-1}\le u_n,\dots,X_{-r}\le u_n \mid X_0>u_n\right).
\]
If $r_2>r_1$, then the event
\[
\{X_{-1}\le u_n,\dots,X_{-r_2}\le u_n\}
\]
is a subset of
\[
\{X_{-1}\le u_n,\dots,X_{-r_1}\le u_n\}.
\]
Conditional probabilities therefore satisfy $\theta_n^{(r_2)} \le \theta_n^{(r_1)}$. Since $\Lambda(u_n,r)=p(u_n)/\theta_n^{(r)}$ with fixed numerator, the second claim follows. The interpretation is immediate: increasing the run length makes the clustering rule stricter, which can only weakly increase the cluster-adjusted stress load.

\begin{proposition}[Finite-sample perturbation bound]
Let $p>0$ and $\theta>0$, and let $\widetilde{p}$ and $\widetilde{\theta}$ be estimators such that
\[
|\widetilde{p}-p|\le \varepsilon_p,
\qquad
|\widetilde{\theta}-\theta|\le \varepsilon_\theta<\theta.
\]
Then
\[
\left|\frac{\widetilde{p}}{\widetilde{\theta}}-\frac{p}{\theta}\right|
\le
\frac{\varepsilon_p}{\theta-\varepsilon_\theta}
\frac{p\,\varepsilon_\theta}{\theta(\theta-\varepsilon_\theta)}.
\]
\end{proposition}

\textit{Proof sketch.}
Write
\[
\frac{\widetilde{p}}{\widetilde{\theta}}-\frac{p}{\theta}
=
\frac{\theta(\widetilde{p}-p)+p(\theta-\widetilde{\theta})}{\theta\widetilde{\theta}}.
\]
Taking absolute values and using $\widetilde{\theta}\ge \theta-\varepsilon_\theta>0$ gives
\[
\left|\frac{\widetilde{p}}{\widetilde{\theta}}-\frac{p}{\theta}\right|
\le
\frac{\theta\varepsilon_p+p\varepsilon_\theta}{\theta(\theta-\varepsilon_\theta)},
\]
which is the stated bound after separating terms. The interpretation is important: estimation error in the clustering term is amplified when $\theta$ is small, so cluster-adjusted stress is most numerically delicate precisely in the regimes where clustering is strongest.

\begin{remark}
The functional $\Lambda(u_n,r)$ is best understood as a ranking and summary device, not as a welfare object or a structural sufficient statistic. Its value lies in comparative stress characterization.
\end{remark}

\begin{remark}
The perturbation bound explains why threshold sensitivity cannot be treated as a cosmetic appendix item. If threshold changes move $\widehat{\theta}_n^{(r)}$ toward zero, the same sampling error in the clustering estimator induces a larger error in $\widehat{\Lambda}_n$. This is a concrete mathematical reason to report threshold and run-length robustness.
\end{remark}

\section{Synthetic Evidence}

\subsection{Core numerical outputs}

Table~\ref{tab:summary} collects the generated quantities used in the analysis.

\begin{table}[H]
\centering
\caption{Core artifact summaries}
\label{tab:summary}
\begin{tabular}{>{\raggedright\arraybackslash}p{0.29\textwidth} >{\raggedright\arraybackslash}p{0.62\textwidth}}
\toprule
Object & Numerical summary \\
\midrule
Extremal index by series & Clustered synthetic stress: $0.673$; isolated synthetic stress: $0.780$; logistic benchmark observable: $1.000$ \\
Return-time distribution & $249$ return times; mean $20.1$ periods; median $14.0$; 90th percentile $40.2$; maximum $90$ \\
Multivariate stress timeline & $120$ dated observations; $12$ joint exceedance dates; maximum of $2$ active series; mean stress score $0.150$ \\
Volatility comparison & Mean rolling volatility during exceedances $0.3429$; outside exceedances $0.2671$; exceedance count $90$ \\
\bottomrule
\end{tabular}
\end{table}

These quantities are not merely descriptive adornments for the manuscript. They delimit the paper's evidentiary scope. Every substantive claim below is tied to these generated outputs or to the corresponding saved source tables.

\subsection{Synthetic clustering benchmarks}

The most direct evidence for the paper's thesis is the extremal-index comparison between the two synthetic stress benchmarks. The clustered benchmark yields an estimated extremal index of $0.673$, whereas the more isolated benchmark yields $0.780$. The gap is moderate rather than dramatic, and the correct interpretation should remain moderate as well. The point is not that the two benchmarks live in wholly different stochastic universes. The point is that recurrence diagnostics detect a difference in temporal organization that would be inadequately summarized by generic language about turbulence or volatility.

This is precisely the setting in which the cluster-adjusted functional becomes interpretable. For equal exceedance incidence, the series with smaller $\theta$ receives a larger stress load. Even when exceedance counts are similar, the clustered benchmark should be treated as carrying greater persistence of stress. This is the operational meaning of the theory section.

\figasset{figures/04_extremal_index_by_series.png}{Estimated extremal indices across benchmark series. Lower values indicate stronger local clustering of exceedances. The result is finite-sample and threshold-dependent, but it recovers the intended ranking between clustered and more isolated stress processes.}

\subsection{Return-time dispersion}

The return-time artifact contains $249$ inter-exceedance intervals with mean $20.1$, median $14.0$, 90th percentile $40.2$, and maximum $90$. This dispersion matters. A stress process can have the same rough event frequency as another process while differing sharply in the distribution of quiet periods between stress recurrences. That information is lost if the analysis stops at exceedance counts or volatility spikes.

The numerical asymmetry between mean and median also suggests that recurrence is not well described by a single ``typical'' waiting time. For stress monitoring, this means that persistence should not be reduced to an average episode length or a binary crisis label. Recurrence structure is a distributional object.

\figasset{figures/03_return_time_distribution.png}{Empirical distribution of return times between exceedances. The wide spread confirms that recurrence is a distributional feature, not a single-duration label.}

\subsection{Multivariate stress activation}

The multivariate illustration records $120$ observations and $12$ joint exceedance dates, with at most two simultaneous active series and mean stress score $0.150$. The correct claim here is descriptive. Stress is sometimes isolated and sometimes joint. That distinction should be measured before any attempt is made to infer contagion or structural propagation. The stress-state index is therefore a screening tool for joint activation, not a substitute for a structural multivariate model.

\figasset{figures/05_multivariate_stress_timeline.png}{Multivariate stress-state timeline. The figure does not identify transmission channels; it records when stress is isolated and when it is joint across observables.}

\subsection{Why the volatility baseline remains necessary}

The repository's baseline comparison shows higher mean rolling volatility during exceedance periods ($0.3429$) than outside them ($0.2671$). That is substantively useful and methodologically expected. The recurrence framework is not introduced to discredit volatility diagnostics. It is introduced because higher volatility does not by itself tell us whether extremes are isolated or clustered.

The right conclusion is complementarity. Volatility helps identify changing scale conditions. Recurrence diagnostics help identify temporal organization in the tail. A serious stress workflow should retain both objects.

\figasset{figures/06_econometric_vs_recurrence_diagnostics.png}{Baseline econometric and recurrence-oriented diagnostics are complementary. Elevated volatility during exceedance periods is informative, but it does not encode clustering intensity or return-time dispersion.}

\section{Limitations and Scope}

The paper has three substantive limits.

First, the theoretical contribution is intentionally narrow. The paper introduces a new functional and proves basic properties for it, but it does not derive new asymptotic theory for extremal-index estimators, optimize threshold selection, or establish finite-sample risk bounds.

Second, the evidence is primarily synthetic. That is appropriate for identifying what the diagnostics measure, but it is not enough to support broad claims about real macro-financial systems. A stronger empirical paper would require defended threshold rules, uncertainty quantification, and domain-specific interpretation.

Third, the current multivariate construction is descriptive. It quantifies simultaneous stress activation; it does not model contagion, network spillovers, or structural transmission.

These limits are not defects to be hidden by tone. They define the paper's correct scope: a methods and theory bridge for recurrence diagnostics in econometric stress analysis.

\section{Conclusion}

Stress analysis should distinguish tail magnitude from tail organization. The first object is well served by the econometric and risk-measure literature. The second requires explicit recurrence diagnostics. This paper developed a reproducible framework for that purpose and introduced a simple cluster-adjusted stress functional that turns tail clustering into an interpretable scalar summary.

The contribution is deliberately bounded. The paper does not claim a new general theory of dependent extremes. It does claim that recurrence can be formalized, estimated, summarized, and integrated with baseline econometric diagnostics in a way that is mathematically explicit and operationally reproducible. The generated synthetic evidence supports exactly that claim and no more.

{\raggedright
\begin{thebibliography}{99}

\bibitem{adrian2016}
Tobias Adrian and Markus K. Brunnermeier.
CoVaR.
\textit{American Economic Review}, 106(7):1705--1741, 2016.

\bibitem{acharya2017}
Viral V. Acharya, Lasse H. Pedersen, Thomas Philippon, and Matthew Richardson.
Measuring systemic risk.
\textit{Review of Financial Studies}, 30(1):2--47, 2017.

\bibitem{bisias2012}
Dimitrios Bisias, Mark Flood, Andrew W. Lo, and Stavros Valavanis.
A survey of systemic risk analytics.
\textit{Annual Review of Financial Economics}, 4:255--296, 2012.

\bibitem{bollerslev1986}
Tim Bollerslev.
Generalized autoregressive conditional heteroskedasticity.
\textit{Journal of Econometrics}, 31(3):307--327, 1986.

\bibitem{brownlees2017}
Christian Brownlees and Robert F. Engle.
SRISK: A conditional capital shortfall measure of systemic risk.
\textit{Review of Financial Studies}, 30(1):48--79, 2017.

\bibitem{coles2001}
Stuart Coles.
\textit{An Introduction to Statistical Modeling of Extreme Values}.
Springer, London, 2001.

\bibitem{embrechts1997}
Paul Embrechts, Claudia Kl\"uppelberg, and Thomas Mikosch.
\textit{Modelling Extremal Events for Insurance and Finance}.
Springer, Berlin, 1997.

\bibitem{engle1982}
Robert F. Engle.
Autoregressive conditional heteroscedasticity with estimates of the variance of United Kingdom inflation.
\textit{Econometrica}, 50(4):987--1007, 1982.

\bibitem{ferro2003}
C. A. T. Ferro and Johan Segers.
Inference for clusters of extreme values.
\textit{Journal of the Royal Statistical Society: Series B}, 65(2):545--556, 2003.

\bibitem{freitas2010}
Ana C. M. Freitas, Jorge M. Freitas, and Mike Todd.
Hitting time statistics and extreme value theory.
\textit{Probability Theory and Related Fields}, 147:675--710, 2010.

\bibitem{hirata1999}
Masayuki Hirata, Beno\^\i t Saussol, and Sandro Vaienti.
Statistics of return times: a general framework and new applications.
\textit{Communications in Mathematical Physics}, 206:33--55, 1999.

\bibitem{leadbetter1983}
M. R. Leadbetter, Georg Lindgren, and Holger Rootz\'en.
\textit{Extremes and Related Properties of Random Sequences and Processes}.
Springer, New York, 1983.

\bibitem{lucarini2016}
Valerio Lucarini, Davide Faranda, Ana C. M. Freitas, Jorge M. Freitas, Mark Holland, Tobias Kuna, Matteo Nicol, Mike Todd, and Sandro Vaienti.
\textit{Extremes and Recurrence in Dynamical Systems}.
Wiley, Hoboken, 2016.

\bibitem{obrien1987}
George L. O'Brien.
Extreme values for stationary and Markov sequences.
\textit{The Annals of Probability}, 15(1):281--291, 1987.

\bibitem{suveges2007}
M\'aria S\"uveges.
Likelihood estimation of the extremal index.
\textit{Extremes}, 10:41--55, 2007.

\end{thebibliography}
}

\end{document}
